% GENERATED FILE. Edit canonical lessons, then run npm run generate:books.
% canonical-source-hash: fnv1a64:929ad233eed10520
% canonical-lessons: FR-C397-fichier, FR-C397-supprimer, FR-C397-copier, FR-C397-coller, FR-C397-sauvegarder, FR-R397-revision-390-397

\chapter{Les fichiers}
\label{ch:fr-les-fichiers}
\hlchaptermodality{397}

\begin{chapteropening}
\textbf{By the end of this chapter:} \emph{I can work with files: delete, copy and paste them, and save my work.}

Complete the last lesson of chapter 397: I can work with files: delete, copy and paste them, and save my work.
\end{chapteropening}

\section[le fichier]{le fichier — the file (on a computer)}

\label{lesson:FR-C397-fichier}



\begin{quote}
Before the new one: say the French for the software, then the French for the update.
\end{quote}

\subsection*{You'll want to know: le fichier}
\textbf{le fichier} — \textquotedblleft{}the file (on a computer)\textquotedblright{}.

\begin{quote}
Je t'ai envoyé le fichier par courriel.
\end{quote}

\begin{quote}
\emph{I sent you the file by email.}
\end{quote}

A paper file is \emph{un dossier}.

\begin{cousinweb}
\emph{fichier} is built on \emph{fiche} \textquotedblleft{}an index card\textquotedblright{}, from \emph{ficher} \textquotedblleft{}to fasten, to fix in\textquotedblright{} (EtymDB), usually traced to a Vulgar Latin \emph{fīgicāre}, from \emph{fīgere} \textquotedblleft{}to fix\textquotedblright{} (Lewis \& Short).
\end{cousinweb}

\begin{grammarlens}[title={What you've built}]
One more word for telling, explaining and describing at B1.
\end{grammarlens}

\subsection*{Your turn}
\begin{itemize}
\raggedright
  \item \emph{Say it:} \emph{le fichier}
  \item \emph{Say it:} \emph{le fichier}, once more
  \item \emph{Say it:} a sentence of your own with \emph{le fichier}
  \item \emph{Recall:} say the French for the software, then the French for the update, then say \emph{le fichier} again
\end{itemize}

\begin{tcolorbox}[breakable,colback=teal!4,colframe=teal!35!black,title={Before you move on}]
What does le fichier mean? (\textquotedblleft{}The file (on a computer)\textquotedblright{}.) Say it once more.
\end{tcolorbox}
\section[supprimer]{supprimer — to delete; to remove}

\label{lesson:FR-C397-supprimer}



\begin{quote}
Before the new one: say the French for the update, then the French for the file.
\end{quote}

\subsection*{You'll want to know: supprimer}
\textbf{supprimer} — \textquotedblleft{}to delete; to remove\textquotedblright{}.

\begin{quote}
J'ai supprimé toutes les vieilles photos.
\end{quote}

\begin{quote}
\emph{I deleted all the old photos.}
\end{quote}

Also to abolish: \emph{supprimer une loi}.

\begin{cousinweb}
\emph{supprimer} is a learned borrowing from Latin \emph{supprimere} \textquotedblleft{}to press under, to hold back\textquotedblright{} (EtymDB; Lewis \& Short), built on \emph{premere} \textquotedblleft{}to press\textquotedblright{}.
\end{cousinweb}

\begin{grammarlens}[title={What you've built}]
One more word for telling, explaining and describing at B1.
\end{grammarlens}

\subsection*{Your turn}
\begin{itemize}
\raggedright
  \item \emph{Say it:} \emph{supprimer}
  \item \emph{Say it:} \emph{supprimer}, once more
  \item \emph{Say it:} a sentence of your own with \emph{supprimer}
  \item \emph{Recall:} say the French for the update, then the French for the file, then say \emph{supprimer} again
\end{itemize}

\begin{tcolorbox}[breakable,colback=teal!4,colframe=teal!35!black,title={Before you move on}]
What does supprimer mean? (\textquotedblleft{}To delete; to remove\textquotedblright{}.) Say it once more.
\end{tcolorbox}
\section[copier]{copier — to copy}

\label{lesson:FR-C397-copier}



\begin{quote}
Before the new one: say the French for the file, then the French for to delete.
\end{quote}

\subsection*{You'll want to know: copier}
\textbf{copier} — \textquotedblleft{}to copy\textquotedblright{}.

\begin{quote}
Il a copié le texte dans un nouveau document.
\end{quote}

\begin{quote}
\emph{He copied the text into a new document.}
\end{quote}

At school, \emph{copier} also means to cheat by copying.

\begin{cousinweb}
\emph{copier} comes from Medieval Latin \emph{copiāre} \textquotedblleft{}to transcribe\textquotedblright{} (EtymDB), from \emph{cōpia} \textquotedblleft{}plenty, abundance\textquotedblright{} (Lewis \& Short): to make more of a text.
\end{cousinweb}

\begin{grammarlens}[title={What you've built}]
One more word for telling, explaining and describing at B1.
\end{grammarlens}

\subsection*{Your turn}
\begin{itemize}
\raggedright
  \item \emph{Say it:} \emph{copier}
  \item \emph{Say it:} \emph{copier}, once more
  \item \emph{Say it:} a sentence of your own with \emph{copier}
  \item \emph{Recall:} say the French for the file, then the French for to delete, then say \emph{copier} again
\end{itemize}

\begin{tcolorbox}[breakable,colback=teal!4,colframe=teal!35!black,title={Before you move on}]
What does copier mean? (\textquotedblleft{}To copy\textquotedblright{}.) Say it once more.
\end{tcolorbox}
\section[coller]{coller — to paste; to stick}

\label{lesson:FR-C397-coller}



\begin{quote}
Before the new one: say the French for to delete, then the French for to copy.
\end{quote}

\subsection*{You'll want to know: coller}
\textbf{coller} — \textquotedblleft{}to paste; to stick\textquotedblright{}.

\begin{quote}
Elle a collé le texte dans son message.
\end{quote}

\begin{quote}
\emph{She pasted the text into her message.}
\end{quote}

\emph{Copier-coller} — copy and paste. Informal: \emph{être collé à un examen} — to fail it.

\begin{cousinweb}
\emph{coller} is built on \emph{colle} \textquotedblleft{}glue\textquotedblright{}, which comes through Vulgar Latin from Greek \emph{kólla} \textquotedblleft{}glue\textquotedblright{} (EtymDB).
\end{cousinweb}

\begin{grammarlens}[title={What you've built}]
One more word for telling, explaining and describing at B1.
\end{grammarlens}

\subsection*{Your turn}
\begin{itemize}
\raggedright
  \item \emph{Say it:} \emph{coller}
  \item \emph{Say it:} \emph{coller}, once more
  \item \emph{Say it:} a sentence of your own with \emph{coller}
  \item \emph{Recall:} say the French for to delete, then the French for to copy, then say \emph{coller} again
\end{itemize}

\begin{tcolorbox}[breakable,colback=teal!4,colframe=teal!35!black,title={Before you move on}]
What does coller mean? (\textquotedblleft{}To paste; to stick\textquotedblright{}.) Say it once more.
\end{tcolorbox}
\section[sauvegarder]{sauvegarder — to save (a file); to safeguard}

\label{lesson:FR-C397-sauvegarder}



\begin{quote}
Before the new one: say the French for to copy, then the French for to paste.
\end{quote}

\subsection*{You'll want to know: sauvegarder}
\textbf{sauvegarder} — \textquotedblleft{}to save (a file); to safeguard\textquotedblright{}.

\begin{quote}
Il a oublié de sauvegarder son travail.
\end{quote}

\begin{quote}
\emph{He forgot to save his work.}
\end{quote}

\emph{Une sauvegarde} — a backup.

\begin{cousinweb}
\emph{sauvegarder} is built from \emph{sauve} \textquotedblleft{}safe\textquotedblright{}, from Latin \emph{salvus} \textquotedblleft{}safe\textquotedblright{} (Lewis \& Short), and \emph{garder} \textquotedblleft{}to keep\textquotedblright{}, usually traced to a Frankish word.
\end{cousinweb}

\begin{grammarlens}[title={What you've built}]
One more word for telling, explaining and describing at B1.
\end{grammarlens}

\subsection*{Your turn}
\begin{itemize}
\raggedright
  \item \emph{Say it:} \emph{sauvegarder}
  \item \emph{Say it:} \emph{sauvegarder}, once more
  \item \emph{Say it:} a sentence of your own with \emph{sauvegarder}
  \item \emph{Recall:} say the French for to copy, then the French for to paste, then say \emph{sauvegarder} again
\end{itemize}

\begin{tcolorbox}[breakable,colback=teal!4,colframe=teal!35!black,title={Before you move on}]
What does sauvegarder mean? (\textquotedblleft{}To save (a file); to safeguard\textquotedblright{}.) Say it once more.
\end{tcolorbox}
\section[(dialogue)]{Retour sur les chapitres 390–397}

\label{lesson:FR-R397-revision-390-397}



\begin{quote}
Say the French for the budget and the income.
\end{quote}

\subsection*{Your turn}
Say these aloud:

\begin{itemize}
\raggedright
  \item the budget, the income, the subscription, the expense, the contribution
  \item the savings, to save up, to invest, the investment, the rate
  \item the cash, the cheque, the chequebook, the bank transfer, to take out
  \item the guarantee, to guarantee, the quote, profitable, good value
  \item the profit, the loss, the bankruptcy, to pay back, the loan
  \item the user, the password, the username, to connect, to access
  \item to install, to uninstall, the device, the software, the update
  \item the file, to delete, to copy, to paste, to save
\end{itemize}

\begin{tcolorbox}[breakable,colback=teal!4,colframe=teal!35!black,title={Before you move on}]
The budget? (\textbf{le budget}.) The income? (\textbf{le revenu}.) The subscription? (\textbf{l'abonnement}.) The expense? (\textbf{la dépense}.) The contribution? (\textbf{la cotisation}.) The savings? (\textbf{l'épargne}.) To save up? (\textbf{épargner}.) To invest? (\textbf{investir}.) The investment? (\textbf{l'investissement}.) The rate? (\textbf{le taux}.) The cash? (\textbf{les espèces}.) The cheque? (\textbf{le chèque}.) The chequebook? (\textbf{le chéquier}.) The bank transfer? (\textbf{le virement}.) To take out? (\textbf{retirer}.) The guarantee? (\textbf{la garantie}.) To guarantee? (\textbf{garantir}.) The quote? (\textbf{le devis}.) Profitable? (\textbf{rentable}.) Good value? (\textbf{avantageux}.) The profit? (\textbf{le bénéfice}.) The loss? (\textbf{la perte}.) The bankruptcy? (\textbf{la faillite}.) To pay back? (\textbf{rembourser}.) The loan? (\textbf{l'emprunt}.) The user? (\textbf{l'utilisateur}.) The password? (\textbf{le mot de passe}.) The username? (\textbf{l'identifiant}.) To connect? (\textbf{connecter}.) To access? (\textbf{accéder}.) To install? (\textbf{installer}.) To uninstall? (\textbf{désinstaller}.) The device? (\textbf{l'appareil}.) The software? (\textbf{le logiciel}.) The update? (\textbf{la mise à jour}.) The file? (\textbf{le fichier}.) To delete? (\textbf{supprimer}.) To copy? (\textbf{copier}.) To paste? (\textbf{coller}.) To save? (\textbf{sauvegarder}.)
\end{tcolorbox}
